\documentclass{beamer}
\usetheme{Madrid}
\title{Nano-Interface-Ready Neuro-Semantic State Communication}
\author{Fernando May Fuentes\\IPN UPIITA / Beihang University\\ORCID: 0009-0002-3953-5224}
\institute{NanoX 2026}
\begin{document}
\frame{\titlepage}
\begin{frame}{Positioning}NDAN is a research architecture for probabilistic cognitive-state interfaces; it is not a claim of direct telepathy or a clinical device.\end{frame}
\begin{frame}{New contribution}Beyond WITCOM 2025 NDAN: typed cognitive-state messages with uncertainty, consent scope, provenance, and revocation.\end{frame}
\begin{frame}{Architecture}Sensing $\rightarrow$ state estimation $\rightarrow$ NSSC message $\rightarrow$ adaptive networking $\rightarrow$ governance.\end{frame}
\begin{frame}{Nanotechnology connection}Future nanosensors, flexible electrodes, low-power edge electronics, and biocompatible packaging are defined as validation targets.\end{frame}
\begin{frame}{Validation}Public BCI data, subject-wise evaluation, uncertainty calibration, privacy, latency, energy, and human-in-the-loop safety.\end{frame}
\end{document}
